\section{Τρέχουσα Κατάσταση και Τάσεις στο \ENG{E-Learning}}

Παρά την εμφάνιση νεότερων τεχνολογικών προσεγγίσεων, το οικοσύστημα του \ENG{E-Learning} εξακολουθεί να
χαρακτηρίζεται από την εκτεταμένη χρήση του προτύπου \ENG{SCORM}. Σε πολλούς οργανισμούς, το βασικό τεχνικό
ζητούμενο δεν είναι η αντικατάσταση του προτύπου, αλλά η ασφαλής και αξιόπιστη αξιοποίηση του ήδη υπάρχοντος
εκπαιδευτικού περιεχομένου.

Παράλληλα έχουν εμφανιστεί νεότερα πρότυπα, όπως το \ENG{xAPI} και το συγγενές \ENG{cmi5}, τα οποία επιχειρούν να
επεκτείνουν τη δυνατότητα καταγραφής δραστηριοτήτων πέρα από το παραδοσιακό μοντέλο των
\ENG{Learning Management Systems (LMS)}. Ωστόσο, η υιοθέτησή τους παραμένει περιορισμένη σε σχέση με τη μεγάλη
εγκατεστημένη βάση \ENG{SCORM} υλικού.

Ταυτόχρονα, η αρχιτεκτονική των συστημάτων μάθησης μετατοπίζεται από μονολιθικές υλοποιήσεις προς πιο αρθρωτά
μοντέλα ή/και \ENG{microservices} - δεδομένου του όγκου των δεδομένων και των ταυτόχρονων συνδέσεων. 
Η εκτέλεση περιεχομένου, η αποθήκευση κατάστασης και η διαχείριση μαθημάτων υλοποιούνται όλο και
συχνότερα ως ανεξάρτητες υπηρεσίες που επικοινωνούν μέσω \ENG{web APIs} και λειτουργούν σε \ENG{cloud} υποδομές.

Σημαντικός παράγοντας σχεδιασμού αποτελεί επίσης η ασφάλεια και το \ENG{monitoring} των συστημάτων. Η χρήση
μηχανισμών αυθεντικοποίησης, η προστασία της πρόσβασης σε εκπαιδευτικό περιεχόμενο και η δυνατότητα
παρακολούθησης της λειτουργίας μιας πλατφόρμας αποτελούν πλέον βασικές απαιτήσεις για σύγχρονες λύσεις
\ENG{E-Learning} και όχι μόνο.